\section{Evaluation}  %3 pages
\label{sec:evaluation}

{\color{blue}

\subsection{Experiment Setup} 
\label{sec:setup}

\subsubsection{Hardware and Software Platforms} 

We evaluate MAGSAT on two GPU platforms: the NVIDIA A100 (Ampere architecture) and the H100 (Hopper architecture). 
These two platforms represent recent NVIDIA GPU generations with distinct differences in hardware resources, 
enabling us to assess whether the proposed framework remains effective across hardware generations. 
All experiments are conducted on machines running Ubuntu 22.04, with CUDA 12.6 and PyTorch 2.9.0.

\subsubsection{Comparison Configurations and Methods}

We evaluate both operator-level and model-level workloads. At the
operator level, we use representative GEMM and convolution workloads,
which cover the compute-intensive primitives commonly appearing in
Transformer and CNN models. At the model level, we evaluate four
representative neural networks: BERT~\cite{devlin2019bert},
LLaMA~\cite{touvron2023llama}, ResNet~\cite{he2016resnet}, and
AlexNet~\cite{krizhevsky2012alexnet}. These models exercise both
GEMM-dominated and Convolution-dominated computation patterns.

We compare MAGSAT with five baselines. \textbf{Torch} denotes eager-mode
PyTorch execution. \textbf{\texttt{torch.compile}} denotes the PyTorch 2.x
compiler stack based on TorchDynamo and TorchInductor~\cite{ansel2024pytorch}.
\textbf{CUTLASS} is NVIDIA's template-based GPU kernel library~\cite{nvidia22cutlass}.
\textbf{DietCode}~\cite{zheng2022dietcode} and
\textbf{Helix}~\cite{zhou2025helix} are state-of-the-art systems in dynamic-shape 
tensor program compilation. All experiments use FP16 inputs and outputs. For timing, we
perform 100 warm-up iterations before measurement to to record the average performance.



\subsection{Operator Performance}
\label{sec:operatorperf}

% 融合算子（单算子）的性能，如果单算子性能一般，就不用放

\subsection{End-to-end Performance}
\label{sec:end2end}

% 端到端网络推理的性能，LLM*2，CNN*2

\subsection{Ablation Study}
\label{sec:ablation}

% 组件的ablation测试，还需要再想想有啥能拆出来的；或者有什么框架特定的设死的参数，可做参数敏感性分析

\subsection{Parameter Tuning Efficiency}
\label{sec:tuning}

% 选到的配置相比全局最优配置的性能差别，以及在各种形状下取到最优配置的比例



\subsection{Overhead Analysis}
\label{sec:overhead}

% 主要是融合方案判定和参数调优的开销比例？

\subsection{Case Study: Qwen3-xxB Inference}
\label{sec:case}

% 可选（时间够用就加）：给个pretrained模型的端到端推理加速会更有说服力

}
